\chapter{Real-Time Messaging over WebSocket}
\label{chap:websocket}

The live path is \verb|GET /ws/chat?token=<jwt>|, handled by
\verb|src/routes/ws.rs|, backed by the connection registry in \verb|src/ws.rs|.
REST endpoints persist and read history; the WebSocket is how a message reaches a
recipient \emph{now}.

\section{The Connection Registry}

\verb|WsConnections| (\verb|src/ws.rs|) is the in-process map of who is online and
how to reach them:

\begin{lstlisting}
pub struct WsConnections {
    inner: Mutex<HashMap<String, Vec<(Uuid, ConnTx)>>>,
}
pub type ConnTx = UnboundedSender<Message>;
\end{lstlisting}

A user maps to a \emph{list} of connections because one account can be signed in
from several devices/tabs at once. Each entry pairs a per-connection \verb|Uuid|
with an unbounded mpsc sender that feeds that socket's writer task. The API:

\begin{center}
\begin{tabular}{@{}lp{0.62\textwidth}@{}}
\toprule
\textbf{Method} & \textbf{Effect} \\
\midrule
\verb|register(user, id, tx)| & Add a connection for \verb|user|. \\
\verb|unregister(user, id)| & Remove it; drop the user key entirely if it was the last connection. \\
\verb|send_to(user, payload)| & Send text to \emph{all} of \verb|user|'s connections. \\
\verb|send_to_except(user, id, payload)| & Send to all but the connection \verb|id| --- used to echo a DM to the sender's \emph{other} devices without looping it back to the origin. \\
\bottomrule
\end{tabular}
\end{center}

Sends are best-effort: \verb|let _ = tx.send(...)|. A dead receiver simply drops
the frame; the connection is cleaned up when its receive loop ends. The registry
has a thorough unit-test module (multi-connection delivery, the ``except''
exclusion, no-op sends to unknown users, and last-connection removal).

\section{Upgrade and the Per-Connection Tasks}

\verb|ws_chat| decodes the token \emph{before} upgrading. An invalid token still
upgrades but is immediately closed with code \verb|4001 "Invalid token"| --- the
handshake completes so the client gets a clean, typed rejection rather than a
bare HTTP error.

On a valid token, \verb|handle_socket| sets up the machinery:

\begin{enumerate}
  \item Mint a \verb|conn_id: Uuid| for this socket.
  \item Split the socket into a \verb|sink| (outbound) and \verb|stream|
        (inbound).
  \item Create an unbounded mpsc channel and spawn a \textbf{writer task} that
        drains the channel into the sink. All outbound frames --- relays, acks,
        errors --- go through this one channel, so any handler (even one running
        for a \emph{different} user) can enqueue a frame for this socket without
        sharing the sink.
  \item \verb|register| the \verb|(conn_id, tx)| under the username.
  \item Loop over inbound frames until the stream ends or a close frame arrives.
  \item On exit, \verb|unregister|, drop the sender, and await the writer task.
\end{enumerate}

\section{Frame Dispatch}

Each inbound text frame is parsed as JSON and dispatched on its \verb|"type"|
field, defaulting to \verb|"dm"|:

\begin{lstlisting}
match data.get("type").and_then(|v| v.as_str()).unwrap_or("dm") {
    "channel" => handle_channel(...).await,
    _         => handle_dm(...).await,
}
\end{lstlisting}

Malformed JSON is logged and skipped, not fatal --- a bad frame never tears down
the connection.

\section{Direct Messages: the Double-Write}

\verb|handle_dm| is where the \verb|channel_id| overloading from
Chapter~\ref{chap:database} pays off. A DM inbound frame looks like:

\begin{lstlisting}
{
  "type": "dm",
  "recipient": "bob",
  "ciphertext": "...", "iv": "...", "sender_public_key_jwk": "...",
  "self_ciphertext": "...", "self_iv": "...", "self_sender_public_key_jwk": "..."
}
\end{lstlisting}

The client encrypts the message \emph{twice}: once to the recipient's public key,
once to its own. The handler therefore writes up to two rows:

\begin{enumerate}
  \item \textbf{The recipient copy} --- \verb|recipient = bob|,
        \verb|channel_id = ''|, using \verb|ciphertext|/\verb|iv|/\verb|spk|.
        This is what Bob reads back via \verb|GET /api/messages/alice|.
  \item \textbf{The self copy} --- only if all three \verb|self_*| fields are
        present --- with \verb|channel_id = 'self'|, using the \verb|self_*|
        values. This is what Alice reads back so she can decrypt her own sent
        history. (Because the recipient copy is encrypted to Bob's key, Alice
        \emph{cannot} decrypt it; she needs her own copy.)
\end{enumerate}

Both rows share the same timestamp. The message id returned to the sender is the
recipient copy's id.

Relay then happens in three directions:

\begin{lstlisting}
state.ws_connections.send_to(recipient, &relay_text);          // to Bob
state.ws_connections.send_to_except(sender, own_id, &relay_text); // Alice's OTHER devices
own_tx.send(ack);                                              // ack to THIS socket
\end{lstlisting}

The relayed frame is a \verb|"dm"| object carrying the id, sender, recipient,
ciphertext, iv, sender public key, timestamp, and message type. The ack is the
same frame plus an \verb|"ack"| field holding the message id.

\section{Channel Messages: Fan-Out}

\verb|handle_channel| parses the frame into the very same
\verb|ChannelMessagePayload| the REST handler uses. It verifies the sender is a
member of \verb|payload.channel_id| (replying with an \verb|{"error": ...}| frame
if not), then delegates to the shared functions in
\verb|src/routes/messages.rs|:

\begin{itemize}
  \item \verb|store_channel_message| --- writes one \verb|messages| row per
        envelope (each encrypted to a different member) plus, for an image, one
        \verb|attachments| row. It returns a \verb|ChannelMessageStoreResult| with
        a \verb|group_id| (shared logical-message id), the timestamp, the
        attachment id, the resolved message type, and the per-envelope row ids.
  \item \verb|relay_channel_message| --- iterates envelopes zipped with their row
        ids and \verb|send_to| each \verb|target_user|, delivering each member
        their own ciphertext.
\end{itemize}

The sender receives an ack keyed by \verb|group_id|. Because storage and relay
are shared functions, a channel message sent over REST and one sent over the
WebSocket are byte-for-byte equivalent in the database and to recipients --- the
transport is an implementation detail.

\section{Why This Shape}

Storing a separate ciphertext per recipient (fan-out) is what keeps the server
zero-knowledge for groups: there is no group key the server could hold. The cost
is storage and write amplification proportional to channel size; the benefit is
that membership changes and per-recipient delivery need no re-keying on the
server, and the server literally cannot read the traffic it carries.
